\documentclass[aip,jmp,preprint,onecolumn,nofootinbib]{revtex4-2}

\usepackage[T1]{fontenc}
\usepackage[utf8]{inputenc}
\usepackage{amsmath,amssymb,amsthm,bm}
\usepackage{mathtools}
\usepackage{microtype}
\usepackage{hyperref}

\hypersetup{
  colorlinks=true,
  linkcolor=blue,
  citecolor=blue,
  urlcolor=blue,
  pdftitle={A hypocycloid spectral bound for circulant unistochastic matrices without symmetric amplitudes},
  pdfauthor={Xiaojian Zeng}
}

\newtheorem{proposition}{Proposition}
\newtheorem{corollary}{Corollary}
\newcommand{\Hd}{\mathcal H_d}
\newcommand{\Cd}{\mathcal C_d}
\newcommand{\Ud}{\mathcal U_d}
\newcommand{\T}{\mathbb T}

\begin{document}

\title{A hypocycloid spectral bound for circulant unistochastic matrices without symmetric amplitudes}

\author{Xiaojian Zeng}
\email{xiaojianzeng@foxmail.com}
\affiliation{PBC School of Finance, Tsinghua University, Beijing 100190, China}

\date{September 13, 2026}

\begin{abstract}
Circulant transition probabilities can arise from a unitary matrix whose
amplitudes do not share the translation symmetry.  We prove that this symmetry
mismatch does not enlarge the possible value of any individual Fourier mode.
For every dimension $d$, if a circulant bistochastic matrix
$B_{ij}=b_{i-j}$ has a same-dimensional unitary realization
$B_{ij}=|U_{ij}|^2$, then each Fourier eigenvalue of $B$ is the normalized
trace of the special-unitary multiplicative commutator
$D_kUD_k^\dagger U^\dagger$.  It therefore lies in the closed unit
$d$-hypocycloid, even when $U$ is not circulant.  At $d=5$ we apply this result
to an exact rational example.  The circulant matrix generated by
$b=(209/500,3/40,7/500,3/40,209/500)$ satisfies every row and column bracelet
inequality, yet its first Fourier eigenvalue lies strictly outside the
five-hypocycloid; hence the matrix is not unistochastic.  The result shows both
that bracelet conditions are insufficient on the five-dimensional circulant
slice and that a marginal hypocycloid violation cannot separate arbitrary
unitary realizations from translation-invariant unitary realizations.  The
corresponding five-site phase-recovery conjecture remains unresolved.
\end{abstract}

\maketitle

\section{Introduction}

A bistochastic matrix describes a finite classical transition law, whereas a
unistochastic matrix has the more restrictive representation
$B_{ij}=|U_{ij}|^2$ for a unitary matrix $U$.  The latter class connects
classical transition probabilities to coherent finite-dimensional dynamics.
Its geometry is substantially less transparent than that of the Birkhoff
polytope, and even useful necessary conditions need not be sufficient.
Rajchel-Mieldzio{\'c} \emph{et al.}\ introduced bracelet matrices as a tractable
superset of the unistochastic matrices and studied in detail the intersection
with circulant matrices.\cite{RajchelMieldzioc2022}

Their analysis distinguishes two notions of symmetry.  A transition matrix can
be circulant even when the unitary that realizes it is not.  If the realizing
unitary is itself circulant, Ref.~\onlinecite{RajchelMieldzioc2022} proves that
every eigenvalue of the transition matrix belongs to the closed unit
$d$-hypocycloid.  This region is also the normalized-trace figure of
$\mathrm{SU}(d)$.\cite{Kaiser2006}  Hypocycloids have appeared more generally
in the spectral study of unistochastic matrices,\cite{Zyczkowski2003} while
block-valued extensions lead to a distinct generalized notion of
unistochasticity.\cite{Nechita2025}

Here we remove the amplitude-level symmetry assumption from the marginal
spectral bound.  The key observation is a commutator-trace identity: on the
circulant slice, every spatial Fourier eigenvalue is the normalized trace of a
multiplicative commutator in $\mathrm{SU}(d)$.  This produces two consequences.
First, a unitary realization whose amplitudes break translation symmetry still
obeys exactly the same one-mode hypocycloid constraint.  Second, a rational
five-dimensional point gives an exact separation between bracelet and
unistochastic matrices on the circulant slice.

The result has a deliberately limited scope.  It does not decide whether every
circulant unistochastic matrix has a circulant unitary realization, a question
conjectured affirmatively in Ref.~\onlinecite{RajchelMieldzioc2022}.  Instead,
it proves that any counterexample must evade detection by each marginal
hypocycloid test and must involve a genuinely joint phase constraint.

\section{Circulant transition laws and trace figures}

Fix an integer $d\geq 2$ and label sites by $\mathbb Z_d$.  Let
$\zeta=\exp(2\pi i/d)$ and define a circulant bistochastic matrix by
\begin{equation}
 B_{ij}=b_{i-j},\qquad b_r\geq 0,\qquad
 \sum_{r=0}^{d-1}b_r=1,
 \label{eq:circulant}
\end{equation}
where all indices are taken modulo $d$.  With the Fourier convention in
Eq.~\eqref{eq:circulant}, its eigenvalues are
\begin{equation}
 \lambda_k(B)=\sum_{r=0}^{d-1}b_r\zeta^{-kr},
 \qquad k=0,\ldots,d-1.
 \label{eq:fourier-eigenvalue}
\end{equation}
The matrix is unistochastic if there exists $U\in \mathrm U(d)$ such that
$B_{ij}=|U_{ij}|^2$ for every pair of sites.  No translation symmetry is
assumed for $U$.

We denote by $\Hd$ the normalized trace figure
\begin{equation}
 \Hd=\left\{\frac{1}{d}\operatorname{tr}W:W\in\mathrm{SU}(d)\right\}.
 \label{eq:trace-figure}
\end{equation}
It is the closed filled unit $d$-hypocycloid.\cite{Kaiser2006}  A boundary
parametrization is
\begin{equation}
 h_d(\theta)=\frac{d-1}{d}e^{i\theta}
              +\frac{1}{d}e^{-i(d-1)\theta}.
 \label{eq:hypocycloid}
\end{equation}

For later use, recall the bracelet condition.  For every pair of rows
$p\neq q$, the $d$ nonnegative lengths
\begin{equation}
 \ell_j=\sqrt{B_{pj}B_{qj}}
 \label{eq:bracelet-lengths}
\end{equation}
must obey the polygon inequality
$2\max_j\ell_j\leq\sum_j\ell_j$, and the analogous condition must hold for
every pair of columns.  These inequalities are necessary for
unistochasticity: they express the possibility that the terms in the
orthogonality relation between two rows or columns close to form a polygon in
the complex plane.\cite{RajchelMieldzioc2022}

\section{A commutator-trace identity}

\begin{proposition}
Let $B$ be a $d\times d$ circulant bistochastic matrix of the form
Eq.~\eqref{eq:circulant}.  If $B_{ij}=|U_{ij}|^2$ for an arbitrary
$U\in\mathrm U(d)$, then, for every $k$,
\begin{equation}
 \lambda_k(B)=\frac{1}{d}\operatorname{tr}
 \left(D_kUD_k^\dagger U^\dagger\right)\in\Hd,
 \qquad
 D_k=\operatorname{diag}\!\left(\zeta^{-ki}\right)_{i=0}^{d-1}.
 \label{eq:main-identity}
\end{equation}
\end{proposition}

\begin{proof}
Set $C_k=D_kUD_k^\dagger U^\dagger$.  This matrix is unitary, and its
determinant cancels multiplicatively:
\begin{equation}
 \det C_k=(\det D_k)(\det U)(\det D_k)^*(\det U)^*=1.
 \label{eq:determinant}
\end{equation}
Thus $C_k\in\mathrm{SU}(d)$.  Expanding its trace and using the fixed
displacement convention gives
\begin{align}
 \frac{1}{d}\operatorname{tr}C_k
 &=\frac{1}{d}\sum_{i,j=0}^{d-1}
     \zeta^{-k(i-j)}|U_{ij}|^2 \notag\\
 &=\frac{1}{d}\sum_{r=0}^{d-1}
     \sum_{i-j=r}\zeta^{-kr}b_r \notag\\
 &=\sum_{r=0}^{d-1}b_r\zeta^{-kr}
 =\lambda_k(B).
 \label{eq:trace-expansion}
\end{align}
There are exactly $d$ ordered pairs $(i,j)$ for each displacement $r$, which
accounts for the cancellation of the prefactor.  Equation
\eqref{eq:trace-figure} now implies $\lambda_k(B)\in\Hd$.
\end{proof}

The proof uses only the symmetry of the squared moduli.  In particular, it
does not assume $[U,S]=0$ for the one-site translation $S$, nor does it change
dimension or introduce an ancilla.  The quantities $\lambda_k(B)$ are the
dimensionless one-step Fourier modes of the site-transition law; the identity
does not imply that two unitary channels agreeing on these site probabilities
agree on coherent inputs.

Relative to the circulant-unitary result of
Ref.~\onlinecite{RajchelMieldzioc2022}, the proposition enlarges the premise
from circulant realizing unitaries to all same-dimensional realizing unitaries,
provided the resulting transition matrix is circulant.  The conclusion remains
a necessary one-mode constraint, not a sufficient condition for either type of
realization.

\section{An exact five-dimensional bracelet witness}

Consider the rational probability vector
\begin{equation}
 b=\left(\frac{209}{500},\frac{3}{40},\frac{7}{500},
          \frac{3}{40},\frac{209}{500}\right)
 \label{eq:witness}
\end{equation}
and let $B_{ij}=b_{i-j}$ for $i,j\in\mathbb Z_5$.

\begin{corollary}
The matrix defined by Eq.~\eqref{eq:witness} is circulant, bistochastic, and
satisfies every row and column bracelet inequality, but it is not
unistochastic.
\end{corollary}

\begin{proof}
Normalization and nonnegativity are immediate.  Write
$p=209/500$, $q=3/40$, and $r=7/500$.  For a row separation of one site, the
bracelet lengths, up to cyclic order, are
\begin{equation}
 \sqrt{pq},\ \sqrt{qr},\ \sqrt{qr},\ \sqrt{pq},\ p.
 \label{eq:lengths-one}
\end{equation}
The only nonautomatic polygon inequality has exact margin
\begin{equation}
 m_1=-\frac{209}{500}+\frac{\sqrt{42}+\sqrt{1254}}{100}.
 \label{eq:margin-one}
\end{equation}
The integer bounds $31350>177^2$ and
$1050>(162/5)^2$ give $m_1>1/1250>0$.  For a separation of two sites, the
length multiset is
\begin{equation}
 \sqrt{pr},\ q,\ \sqrt{pr},\ \sqrt{pq},\ \sqrt{pq}.
 \label{eq:lengths-two}
\end{equation}
Its largest member $\sqrt{pq}$ occurs twice, so the polygon inequality follows
immediately.  Separations three and four reproduce these two multisets, and
the column tests are identical by cyclic reindexing.  Hence all bracelet
conditions hold.

For the first Fourier mode, exact cyclotomic reduction yields
\begin{equation}
 \lambda_1(B)=t e^{i\pi/5},\qquad
 t=\frac{465+343\sqrt{5}}{2000}.
 \label{eq:lambda-witness}
\end{equation}
The radial endpoint of the filled five-hypocycloid on this symmetry ray is
\begin{equation}
 h_5(\pi/5)=\frac{3}{5}e^{i\pi/5}.
 \label{eq:radial-boundary}
\end{equation}
Since $\sqrt5>15/7$, Eq.~\eqref{eq:lambda-witness} gives $t>3/5$ strictly.
Thus $\lambda_1(B)\notin\mathcal H_5$.  The proposition rules out every
same-dimensional unitary realization, proving that $B$ is not unistochastic.
\end{proof}

This is not a counterexample to the phase-recovery conjecture: the witness
fails its essential premise because it has no arbitrary unitary realization.
Instead, it gives a five-dimensional negative answer to the narrower question
whether all bracelet conditions suffice for unistochasticity on the circulant
slice.

\section{Consequence for translation-symmetric realization}

Let $\Cd$ denote the circulant bistochastic matrices, $\Ud$ the unistochastic
matrices, and $\mathcal C_d^2$ the circulant matrices that admit a circulant
unitary realization, following Ref.~\onlinecite{RajchelMieldzioc2022}.  That
work proves
\begin{equation}
 \Cd\cap\Ud=\mathcal C_d^2\qquad (d\leq4)
 \label{eq:small-d}
\end{equation}
and conjectures the equality in every dimension.  The first unsettled case in
that program is
\begin{equation}
 \mathcal C_5\cap\mathcal U_5\stackrel{?}{=}\mathcal C_5^2.
 \label{eq:p5}
\end{equation}

The proposition shows that every element on the left-hand side of
Eq.~\eqref{eq:p5} already satisfies every marginal five-hypocycloid constraint.
The same constraints were previously known for the right-hand side.  Therefore
an eigenvalue outside $\mathcal H_5$ can reject a circulant bistochastic or
bracelet candidate, as in the corollary, but it can never distinguish the two
realization classes in Eq.~\eqref{eq:p5}.  If the equality fails, its separator
must retain information about the joint compatibility of the nontrivial
Fourier modes rather than test them independently.

Physically, the result is a one-step control limitation on a fixed finite ring.
Even when the amplitude matrix breaks translation symmetry, a
translation-invariant table of all site-to-site probabilities cannot push an
individual spatial Fourier mode beyond the same trace figure obeyed by a
translation-invariant amplitude matrix.  This statement is exact and
dimensionless.  It invokes no perturbative order, continuum or thermodynamic
limit, coin degree of freedom, environment, mixture, or postselection.

\section{Methods and reproducibility}

The proof of Eq.~\eqref{eq:main-identity} is an exact index calculation.  The
five-dimensional certificate uses rational arithmetic, symbolic radicals, and
strict integer inequalities; no fitted parameter, optimizer, sampling
argument, or floating-point tolerance enters the conclusion.  A deterministic
SymPy script supplied with the accompanying reproducibility package checks the
normalization, every bracelet class, Eq.~\eqref{eq:lambda-witness}, and the
strict hypocycloid exclusion.  It also tests the sign convention on an exact
unitary whose squared moduli are circulant although the unitary does not commute
with translation.

OpenAI Codex (OpenAI; model reported by the research environment as GPT-5,
session date September 13, 2026) was used to assist with derivation, exact
symbolic checking, bounded literature screening, and manuscript drafting.  Its
outputs are not treated as independent experimental data or external peer
review.  The named author must critically review the calculations, references,
scope, and wording and assume responsibility for any submitted version.

\section{Conclusion}

Every Fourier eigenvalue of a circulant unistochastic matrix is a normalized
special-unitary commutator trace, whether or not the realizing unitary is
circulant.  The resulting all-$d$ hypocycloid bound turns an explicit rational
five-site bracelet matrix into an exact nonunistochastic witness and proves
that marginal spectral exclusion cannot resolve the distinction between
arbitrary and translation-symmetric unitary realizations.  The full five-site
phase-recovery question in Eq.~\eqref{eq:p5} remains unresolved; the next
structural target is the joint compatibility of its nontrivial Fourier modes.

\section*{Author declarations}

\subsection*{Conflict of interest}
The conflict-of-interest declaration must be confirmed by the author before
submission.

\subsection*{Author contributions}
The CRediT contribution statement must be confirmed by the author before
submission.

\subsection*{Data availability}
No empirical data were generated.  The exact symbolic verification script,
software requirements, regeneration instructions, and reference output are
included in the accompanying reproducibility package.

\bibliography{manuscript/references}

\end{document}
